% !TEX root = ../main.tex
\chapter{Objetivos}\label{cap:objetivos}

El estado de la técnica del capítulo anterior sitúa la brecha entre promesas teóricas del QML / optimización cuántico-inspirada y lo medible en régimen NISQ \citep{preskill2018, bowles2024}. Este capítulo fija el objetivo general y los específicos del estudio.

\section{Objetivo general}

El objetivo general de este trabajo es \textbf{evaluar de forma empírica y sistemática} si los algoritmos cuánticos disponibles en la era NISQ ofrecen una \textbf{ventaja real y medible} frente a los algoritmos clásicos consolidados en problemas representativos de la ciencia de datos, con el fin de determinar las condiciones bajo las cuales la computación cuántica puede aportar valor práctico en este dominio \citep{preskill2018, biamonte2017}.

\section{Objetivos específicos}

\begin{enumerate}
    \item \textbf{Definir un marco conceptual} que organice el estudio en torno a \textbf{clasificación}, \textbf{optimización} y \textbf{representación} (feature map / QUBO--Ising) \citep{lloyd2020, pennylaneEmbedBlog, parzygnat2025}.

    \item \textbf{Diseñar e implementar pipelines de clasificación} comparables entre modelos clásicos y cuánticos bajo un \textbf{protocolo justo} de partición, escalado y reducción dimensional cuando el presupuesto de qubits lo exija \citep{schuld2021, havlicek2019, bowles2024}.

    \item \textbf{Diseñar e implementar pipelines de optimización combinatoria} comparables entre solvers clásicos y cuántico-inspirados (incluidos QAOA y VQE) sobre \textbf{formulaciones QUBO/Ising compartidas} \citep{farhi2014, peruzzo2014, lucas2014, lubinski2024}.

    \item \textbf{Evaluar calidad de la solución y esfuerzo computacional} con métricas comunes, de forma \textbf{reproducible} \citep{bowles2024, cerezo2021}.

    \item \textbf{Interpretar los resultados} en el régimen NISQ de pocos qubits y simulación clásica, \textbf{delimitando qué se puede y no se puede afirmar} sobre ventaja cuántica \citep{huang2021, thanasilp2022, mcclean2018, preskill2018}.
\end{enumerate}
